\subsection{Multilinear morphisms}

7.3.1. Let \(M_{1},\ldots,M_{d}\) and N be vector bundles with base B, and let u be a map from the set \(M_{1}\times_{B}\ldots\times_{B}M_{d}\) into N. One says that u is a multilinear morphism (or d-linear) if the following condition is satisfied:

For every \(b_0 \in \mathbf{B}\) , there exists an open neighborhood \(\mathbf{U}\) of \(b_0\) in \(\mathbf{B}\) , vector charts \(t^j = (\mathbf{U}, \varphi^j, \mathbf{F}^j)\) of \(\mathbf{M}_j\) (for \(1 \leqslant j \leqslant d\) ) and \(t = (\mathbf{U}, \varphi, \mathbf{F})\) of \(\mathbf{N}\) , and a map \(\lambda\) of class \(\mathbf{C}^r\) from \(\mathbf{U}\) into the Banach space \(\mathcal{L}(\mathbf{F}^1, \ldots, \mathbf{F}^d; \mathbf{F})\) of continuous d-linear maps from \(\mathbf{F}^1 \times \cdots \times \mathbf{F}^d\) to \(\mathbf{F}\) , such that:

\[
(t _ {b} \circ \lambda (b)) (x _ {1}, \dots , x _ {d}) = u (t _ {b} ^ {1} (x _ {1}), \dots , t _ {b} ^ {d} (x _ {d}))
\]

for all \(b \in \mathbf{U}\) and all \(x_{j} \in \mathbf{F}^{j}\) .

Every multilinear morphism \(u\) is a morphism of manifolds from the fibered product \(\mathbf{M}_{1} \times_{\mathbf{B}} \cdots \times_{\mathbf{B}} \mathbf{M}_{d}\) to \(\mathbf{N}\) and induces for every \(b \in \mathbf{B}\) a continuous \(d\)-linear map \(u_b\) from \((\mathbf{M}_{1})_{b} \times \cdots \times (\mathbf{M}_{d})_{b}\) to \(\mathbf{N}_{b}\) .

If \(f\) is a morphism from B into a manifold \(\mathbf{B}'\) , a multilinear morphism from \(\mathbf{M}_1 \times_{\mathbf{B}} \cdots \times_{\mathbf{B}} \mathbf{M}_d\) into a vector bundle \(\mathbf{M}'\) with base \(\mathbf{B}'\) is by definition the composite of a multilinear morphism from \(\mathbf{M}_1 \times_{\mathbf{B}} \cdots \times_{\mathbf{B}} \mathbf{M}_d\) to \(f^*\mathbf{M}'\) with the canonical \(f\)-morphism of \(f^*\mathbf{M}'\) to \(\mathbf{M}'\) .

A bilinear morphism is also called a pairing. For \(d=1\) , a linear morphism is a morphism in the sense of 7.2.1. For \(d=0\) , a 0-linear morphism is identified with a section of N (7.4).

7.3.2. An algebra bundle over base B is a vector bundle A over base B, equipped with a pairing of \(A \times_{B} A\) into A. Each fiber \(A_{b}\) is then equipped with a K-algebra structure. If for every \(b \in B\) the algebra \(A_{b}\) has an identity element, denoted \(e_{b}\) , the map \(b \mapsto e_{b}\) is a section of A (cf. 7.4). An algebra bundle A is said to be locally trivial if, for every point \(b_{0}\) of B, there exists a vector chart \(t = (\mathrm{U}, \varphi, \mathrm{E})\) from A at the point \(b_{0}\) and a K-algebra structure on E such that \(t_{b}\) is an isomorphism of algebras from E onto \(A_{b}\) for all \(b \in U\) .

7.3.3. Let A be a bundle of associative algebras with unit element over base B. One calls a bundle of A-modules over base B a vector bundle M over base B equipped with a pairing \(m: A \times_{B} M \to M\) such that the map \(m_{b}: A_{b} \times M_{b} \to M_{b}\) defines, for all \(b \in B\) , a structure of \(A_{b}\)-module on the fiber \(M_{b}\) .

Let M and M′ be two vector bundles in A-modules. We call an A-homomorphism from M to M′ any morphism \(g: M \rightarrow M'\) of vector bundles over the base B which induces for every b in B an \(A_b\)-linear map from \(M_b\) to \(M_b'\) .

Suppose A is locally trivial. We say that a bundle M of A-modules is locally trivial if for every point \(b_0\) of B, there exists a vector chart \(t = (U, \varphi, E)\) from A at \(b_0\) as in 7.3.2, a vector chart \(t' = (U, \varphi', L)\) from M at \(b_0\) , and a structure of E-module on L such that, for every \(b \in U\) , \(t_b'\) is a \(t_b\)-isomorphism of the \(A_b\)-module \(M_b\) onto the E-module L.

7.3.4. Let A be a Banach algebra over K (for example, a field equipped with a finite-dimensional K-algebra structure). The trivial bundle \(A_{B}\) is then a locally trivial algebra bundle. A locally trivial bundle M in \(A_{B}\)-modules is also called a vector bundle over A (with base B). The fibers \(M_{b}\) are then topological A-modules. An f-morphism u from M to another vector bundle over A is said to be A-linear if the maps \(u_{b}\) are A-linear for all \(b \in B\) .

